\documentclass[11pt,a4paper]{article}

\usepackage[utf8]{inputenc}

\usepackage[T1]{fontenc}

\usepackage[polish]{babel}

\usepackage{lmodern}

\usepackage{geometry}

\geometry{margin=27mm}

\usepackage[hidelinks]{hyperref}

\title{Długotrwałe używanie agenta jako regulatora emocji a własne strategie samoregulacji}

\author{Mateusz Pomianek}

\date{Szkic naukowy, wersja 1.0.0 | 28 listopada 2026}

\begin{document}

\maketitle

\begin{abstract}

Proponowana miara główna to liczba różnych samodzielnie zastosowanych strategii w kolejnym tygodniu.

\end{abstract}

\section{Wprowadzenie}

Odciążanie poznawcze polega na przenoszeniu części operacji na zasób zewnętrzny (Risko \& Gilbert, 2016). Przeniesienie tego mechanizmu na samoregulację jest propozycją, którą można zestawić z koncepcją elastycznego użycia strategii (Bonanno \& Burton, 2013).

\section{Cel i pytania badawcze}

Hipotezy robocze wymagają zatwierdzenia przed rozpoczęciem właściwej weryfikacji.

\paragraph{H1.} W tygodniach o większym niż indywidualna średnia udziale trudnych sytuacji obsłużonych z pomocą agenta osoba wymieni mniej odrębnych samodzielnie zastosowanych strategii w kolejnym tygodniu, przy tej samej liczbie obserwowanych sytuacji.

\section{Materiał i metody}

Pierwszorzędowy wskaźnik: liczba różnych samodzielnie zastosowanych strategii w kolejnym tygodniu.

Zastosowanie strategii jest kodowane w dzienniczku według zamrożonej rubryki. Model panelowy rozdzieli efekty wewnątrz i między osobami, kontrolując wynik poprzedniego tygodnia i intensywność stresorów. Sprawdzi się także odwrotny kierunek czasowy, ponieważ selekcja użycia agenta ogranicza interpretację przyczynową.

Przed badaniem z udziałem ludzi należy zatwierdzić wymagania etyczne, zasady zgody i ochronę danych.

\section{Plan analizy}

Zamrożenie estymandu, reguł oceny jakości i obsługi braków nastąpi przed dostępem do danych wynikowych. Raport obejmie oszacowanie kontrastu głównego z przedziałem niepewności i pełne mianowniki. Jeśli przewidziano hipotezy drugorzędowe, przed analizą należy ustalić sposób kontroli wielokrotności.

\section{Wyniki}

W niniejszym pliku nie umieszczono wyników empirycznych, statystyk ani nieprzeprowadzonych analiz. Sekcja wymaga uzupełnienia po uzyskaniu i weryfikacji materiału dowodowego.

\section{Dyskusja i ograniczenia}

Ograniczenia obejmują możliwe błędy pomiaru, dobór próby, alternatywne mechanizmy i ograniczoną możliwość uogólniania. Korelacja nie dowodzi przyczynowości; poprawność techniczna nie jest dowodem znaczenia klinicznego.

\section{Wnioski}

Hipotezy i sposób pomiaru stanowią podstawę oceny wykonalności oraz przygotowania właściwego protokołu.

\section*{Bibliografia}

{\setlength{\parindent}{0pt}
\setlength{\parskip}{5pt}
\hangindent=1.27cm
Bonanno, G. A., \& Burton, C. L. (2013). Regulatory flexibility: An individual differences perspective on coping and emotion regulation. \textit{Perspectives on Psychological Science, 8}(6), 591--612. \url{https://doi.org/10.1177/1745691613504116}\par
\hangindent=1.27cm
Risko, E. F., \& Gilbert, S. J. (2016). Cognitive offloading. \textit{Trends in Cognitive Sciences, 20}(9), 676--688. \url{https://doi.org/10.1016/j.tics.2016.07.002}\par
}

\end{document}
